% ============================================================================
% Fichier  : partie6/P06_C16_mesurer_incertitude.tex
% Titre    : Mesurer l'Incertitude d'une Enquête
% Auteur   : MINKA MI NGUIDJOÏ Thierry Emmanuel
% Version  : 1.0 -- novembre 2026
% Corrigés : corriges/Corr_C16.tex
% ============================================================================
% Sources : Minka (2026), A0 « Vers l'auditologie » (HAL hal-05701007) et A1
% sur l'instrument trichotomique (HAL hal-05701270) ; Shannon (1948).
% Les articles A2 à A5 du programme sont en préparation et ne sont pas cités.
% Mesures reprises du corpus (NRAC, ACSI) : mu = part de OUI, U = part de
% NEK, H = entropie de Shannon de la répartition OUI/NON/NEK, en bits.
% Faits nouveaux (bible) : douze propositions d'enquête P1 à P12 cotées au
% 14/03, au 29/03 et au 06/04/2025 (0/0/12 ; 4/2/6 ; 5/2/5).
% ============================================================================

\chapter{Mesurer l'Incertitude d'une Enquête}
\label{chap:auditologie}
\chaptermark{Mesurer l'incertitude}

\epigraph{%
    « Ce que l'on ne sait pas se compte,\\
    comme ce que l'on sait. »%
}{Principe du cours}

\niveauChapitre{Recherche}{%
    Maîtrise de l'ensemble du cours
}

\begin{boxobjectifs}
\begin{itemize}
    \item Représenter l'état d'une enquête par un ensemble de propositions
          cotées \OUI{}, \NON{} ou \NEK{}.
    \item Calculer et interpréter les mesures $\mu$, $U$ et $H$.
    \item Suivre l'évolution de l'incertitude au cours d'une enquête.
    \item Distinguer un \NEK{} réductible d'un \NEK{} irréductible dans le
          cadre d'une mission.
    \item Formuler une question de recherche sur l'incertitude en
          investigation numérique.
\end{itemize}
\end{boxobjectifs}

\bigskip

% ============================================================================
\section*{Avant de commencer : « avez-vous avancé~? »}
% ============================================================================

Bafoussam, 29 mars 2025. Le directeur général de NYANGA CRÉDIT demande à
Mme~MBARGA où en est l'enquête. Elle pourrait répondre par une liste de
découvertes. Elle répond par trois nombres~: sur les douze propositions que
l'enquête doit trancher, quatre sont établies, deux réfutées, six encore
ouvertes. Le 14~mars au soir, les douze étaient ouvertes. Le directeur
comprend ce que l'enquête sait, et ce qu'il lui reste à apprendre.

\begin{boxalert}
\textbf{Violation typique.} Rendre compte de l'avancement d'une enquête par
ce qui a été trouvé, sans dire ce qui reste inconnu ni si cet inconnu
diminue.

\smallskip
\textbf{Règle~:} l'état d'une enquête se décrit par l'ensemble des
propositions qu'elle doit trancher, chacune cotée. L'avancement se mesure à
la baisse de la part des \NEK{}, pas au nombre de découvertes.
\end{boxalert}

\separateur

% ============================================================================
\section{L'Instrument Trichotomique}
\label{sec:c16-instrument}
% ============================================================================

Tout au long du cours, chaque conclusion a reçu une cote \OUI{}, \NON{} ou
\NEK{}. Cet instrument vient du programme de recherche de
l'auditologie~\cite{minka2026auditologie}, qui propose de traiter l'audit
des systèmes d'information, et plus largement toute démarche qui conclut
sur des preuves, comme une science appliquée de la réduction de
l'incertitude. L'article qui formalise l'instrument~\cite{minka2026nek}
montre qu'une réponse binaire, oui ou non, force l'auditeur à transformer
son ignorance en affirmation, et qu'une troisième réponse, \NEK{}, rend
cette ignorance visible et mesurable.

\begin{boxdef}
Soit un ensemble de $n$ propositions, cotées, avec $n_{O}$ réponses \OUI{},
$n_{N}$ réponses \NON{} et $n_{K}$ réponses \NEK{}. On définit~:

\smallskip
\textbf{\termecle{Taux d'établissement}}~: $\mu = n_{O}/n$, la part des
propositions établies.

\smallskip
\textbf{\termecle{Taux d'incertitude}}~: $U = n_{K}/n$, la part des
propositions que les preuves ne permettent pas de trancher.

\smallskip
\textbf{\termecle{Entropie}}~: $H = -\sum_{c} p_{c} \log_{2} p_{c}$, où
$p_{c}$ est la part de chaque cote, exprimée en bits~\cite{shannon1948}.
Elle mesure la dispersion des réponses~: nulle si toutes les cotes sont
identiques, maximale ($\log_{2} 3 \approx 1{,}585$~bit) si elles sont
également réparties.
\end{boxdef}

Dans un audit, \OUI{} signifie qu'un contrôle est en place et \NON{} qu'il
ne l'est pas~; $\mu$ y mesure la conformité. Dans une enquête, \OUI{} et
\NON{} sont deux formes de savoir~: une proposition réfutée est aussi
tranchée qu'une proposition établie. La part de savoir y vaut donc $1 - U$,
et $U$ devient la mesure principale.

% ============================================================================
\section{Les Douze Propositions de NYANGA}
\label{sec:c16-nyanga}
% ============================================================================

\begin{boxscenario}{Opération NYANGA. Trois instantanés de l'enquête}
\begin{tableaucours}{cL{6.2cm}ccc}{%
    \tableauHeader{N°} &
    \tableauHeader{Proposition} &
    \tableauHeader{14/03} &
    \tableauHeader{29/03} &
    \tableauHeader{06/04}
}
P1 & L'attaquant est entré par la passerelle & \NEK{} & \OUI{} & \OUI{} \\
P2 & Par l'exploitation de VULN-2024-4471 & \NEK{} & \NEK{} & \OUI{} \\
P3 & La clé de l'interface a été lue sur SRV-CPT01 & \NEK{} & \OUI{} &
\OUI{} \\
P4 & Les identifiants de \texttt{adm.kamga} ont été extraits de la copie
du 9~mars & \NEK{} & \NEK{} & \NEK{} \\
P5 & Les données sorties sont le dossier \texttt{Clients} & \NEK{} & \OUI{}
& \OUI{} \\
P6 & Les virements ont été émis hors du réseau du siège & \NEK{} & \OUI{} &
\OUI{} \\
P7 & Un employé a participé & \NEK{} & \NEK{} & \NEK{} \\
P8 & L'attaquant a agi du 3 au 8~mars & \NEK{} & \NEK{} & \NEK{} \\
P9 & Les auteurs sont identifiés & \NEK{} & \NEK{} & \NEK{} \\
P10 & La passerelle contient un code persistant & \NEK{} & \NEK{} & \NEK{} \\
P11 & Les virements sont passés par le lien du siège & \NEK{} & \NON{} &
\NON{} \\
P12 & \texttt{upd.exe} était présent sur SRV-FIC01 en 2017 & \NEK{} & \NON{} &
\NON{} \\
\end{tableaucours}
\end{boxscenario}

\begin{boxlecture}{L'incertitude de l'enquête, trois dates}
\begin{tableaucours}{L{2.0cm}cccc}{%
    \tableauHeader{Date} &
    \tableauHeader{$n_{O}/n_{N}/n_{K}$} &
    \tableauHeader{$\mu$} &
    \tableauHeader{$U$} &
    \tableauHeader{$H$ (bit)}
}
14/03 & 0/0/12 & 0 & 1 & 0 \\
29/03 & 4/2/6 & 0,333 & 0,500 & 1,459 \\
06/04 & 5/2/5 & 0,417 & 0,417 & 1,483 \\
\end{tableaucours}

\medskip
\textbf{Calcul pour le 29~mars.} $p_{O} = 4/12$, $p_{N} = 2/12$, $p_{K} =
6/12$~:
\[
H = -\left(\tfrac{1}{3}\log_{2}\tfrac{1}{3} + \tfrac{1}{6}\log_{2}\tfrac{1}{6}
+ \tfrac{1}{2}\log_{2}\tfrac{1}{2}\right) \approx 0{,}528 + 0{,}431 + 0{,}500
\approx 1{,}459~\text{bit}.
\]

\textbf{Lecture.} $U$ passe de 1 à 0,500 puis à 0,417~: l'enquête avance.
$H$, lui, augmente~: il était nul quand tout était inconnu, et la
répartition des cotes devient plus mélangée à mesure que l'enquête tranche.
$H$ ne mesure donc pas l'ignorance, mais la diversité des réponses. La
règle du corpus s'applique~: on lit $U$ d'abord, $\mu$ ensuite, et $H$ pour
contrôler que la répartition n'est pas dégénérée.

\medskip
\hfill\lecture{0{,}417}{0{,}417}{1{,}483}
\end{boxlecture}

% ============================================================================
\section{Deux Sortes de NEK}
\label{sec:c16-nek}
% ============================================================================

Les cinq \NEK{} du 6~avril ne sont pas de même nature. Certains peuvent être
levés par une pièce que l'enquête est en mesure d'obtenir~; d'autres ne le
peuvent pas dans le cadre de la mission.

\begin{tableaucours}{cL{5.2cm}L{2.4cm}L{3.8cm}}{%
    \tableauHeader{N°} &
    \tableauHeader{Ce qui lèverait le \NEK{}} &
    \tableauHeader{Nature} &
    \tableauHeader{Coût ou condition}
}
P4 & Aucune pièce~: l'extraction a eu lieu hors du système & Irréductible &
Hors de portée de toute enquête technique sur le système \\
P7 & Pièces sur les personnes (témoignages, téléphonie) & Réductible hors
mission & Enquête judiciaire \\
P8 & Journaux des agences, s'ils existent encore & Réductible & Quelques
jours d'examen \\
P9 & Données des hébergeurs de 203.0.113.18, .47 et .62 & Réductible hors
mission & Entraide internationale, mois \\
P10 & Image du système de la passerelle par l'éditeur & Réductible &
Réquisition de l'éditeur, semaines \\
\end{tableaucours}
\vspace{-0.8em}
\begin{center}{\small\itshape Tableau~16.1. Les \NEK{} de NYANGA au 6~avril}\end{center}

Cette distinction change la décision. Il est rationnel de lever P8, peu
coûteux. Il est inutile de promettre de lever P4. Lever P9 et P7 appartient
à d'autres acteurs, et le rapport doit le dire pour que la direction ne
l'attende pas de son cabinet.

\begin{boxinfo}
\textbf{La valeur d'une pièce.} Chaque pièce supplémentaire a un coût
(temps, argent, délai) et une valeur~: les \NEK{} qu'elle peut lever, et
l'importance de ces \NEK{} pour la décision. Choisir les pièces à chercher
revient à comparer ces valeurs. Formaliser cette comparaison pour
l'investigation numérique est l'une des questions ouvertes du programme de
recherche.
\end{boxinfo}

% ============================================================================
\section{Questions Ouvertes pour la Recherche}
\label{sec:c16-recherche}
% ============================================================================

L'auditologie a été proposée pour l'audit des systèmes d'information. Son
transfert à l'investigation numérique et à l'analyse de vulnérabilités
soulève des questions qui peuvent fonder un mémoire de master ou une thèse.

\begin{enumerate}
    \item \textbf{La granularité des propositions.} Les mesures dépendent du
          découpage de l'enquête en propositions. Comment définir un
          découpage stable, pour que deux enquêteurs obtiennent des $U$
          comparables sur la même affaire~?
    \item \textbf{La pondération.} Toutes les propositions ne pèsent pas
          autant pour la décision. Faut-il pondérer $U$, et selon quel
          critère qui ne soit pas arbitraire~?
    \item \textbf{La confiance.} Le cours associe à chaque analyse un niveau
          de confiance (faible, moyenne, élevée) distinct de la cote. Comment
          articuler ces deux échelles~?
    \item \textbf{Le programme de vulnérabilités.} Un actif hors inventaire
          est un \NEK{} (Chapitre~\ref{chap:identifier}). Le $U$ d'un parc,
          calculé sur la présence de failles, peut-il servir d'indicateur de
          pilotage, à côté du nombre de failles~?
    \item \textbf{L'attribution.} L'identité des auteurs reste souvent
          \NEK{}. Peut-on décrire formellement les conditions qui la rendent
          réductible~?
\end{enumerate}

\begin{boxinfo}
\textbf{Construire un sujet de mémoire sur l'incertitude}
\begin{enumerate}
    \item Choisir un objet~: une enquête, un programme de vulnérabilités,
          un type de pièce.
    \item Définir l'ensemble des propositions et la règle de cotation.
    \item Recueillir des cas réels anonymisés ou des cas construits,
          documentés.
    \item Mesurer $U$, $\mu$ et $H$, et leur évolution.
    \item Confronter les mesures à un critère extérieur (décision prise,
          erreur constatée a posteriori).
    \item Discuter les limites de la mesure, et proposer une amélioration.
\end{enumerate}
\end{boxinfo}

% ============================================================================
\section{Ce que la Mesure demande à chaque Métier}
\label{sec:c16-metiers}
% ============================================================================

\begin{boxmetier}{Analyste en réponse à incident}
Ouvrir chaque mission par la liste des propositions à trancher, la coter
chaque semaine et rendre compte de l'avancement par $U$. La liste elle-même
oblige à dire, dès le départ, ce que l'on cherche.
\end{boxmetier}

\begin{boxmetier}{Application OPJ/Expert Judiciaire}
Le tableau des propositions cotées, annexé au rapport, donne au tribunal une
vue d'ensemble honnête~: ce qui est établi, ce qui est réfuté, ce qui ne
peut pas l'être.
\end{boxmetier}

\begin{boxmetier}{Analyste en gestion des vulnérabilités}
Publier, à côté du nombre de failles, la part du parc dont on ne sait rien.
C'est le $U$ du programme.
\end{boxmetier}

\begin{boxmetier}{RSSI}
Demander à chaque enquête et à chaque programme ce qu'ils ont trouvé, puis
ce qu'ils ignorent encore et ce qu'il en coûterait de le savoir.
\end{boxmetier}

\separateur

\begin{boxresume}
\textbf{Ce qu'il faut retenir de ce chapitre~:}
\begin{itemize}
    \item L'état d'une enquête se décrit par des propositions cotées
          \OUI{}, \NON{} ou \NEK{}.
    \item $\mu$ est la part des \OUI{}, $U$ la part des \NEK{}, $H$
          l'entropie de la répartition~; dans une enquête, $U$ est la mesure
          principale.
    \item À NYANGA, $U$ passe de 1 à 0,417 en trois semaines, tandis que
          $H$ augmente~: $H$ mesure la diversité, pas l'ignorance.
    \item Un \NEK{} est réductible ou non dans le cadre d'une mission~; la
          distinction guide le choix des pièces à chercher.
    \item Le transfert de l'auditologie à l'investigation ouvre des
          questions de recherche sur la granularité, la pondération et la
          confiance.
\end{itemize}
\end{boxresume}

\begin{boxinfo}
\textbf{Points de contrôle.} Vous avez compris ce chapitre si vous savez~:
\begin{itemize}
    \item calculer $\mu$, $U$ et $H$ pour un ensemble de cotes~;
    \item expliquer pourquoi $H$ peut augmenter pendant qu'une enquête
          avance~;
    \item classer un \NEK{} comme réductible ou irréductible.
\end{itemize}
\end{boxinfo}

% ============================================================================
\section*{Exercices}
% ============================================================================

\begin{boxexo}[title={\ding{45}\quad Exercice~16.1. Calculer \niveaubase}]
Calculez $\mu$, $U$ et $H$ pour les répartitions suivantes de dix
propositions~: (a) 10/0/0~; (b) 5/5/0~; (c) 4/3/3~; (d) 0/0/10.
\end{boxexo}

\begin{boxexo}[title={\ding{45}\quad Exercice~16.2. Le 6 avril \niveaumoyen}]
\begin{enumerate}[(a)]
    \item Vérifiez le calcul de $H$ au 6~avril.
    \item Si l'examen des journaux des agences établit l'activité du 3 au
          8~mars (P8 passe à \OUI{}), que deviennent $\mu$, $U$ et $H$~?
    \item Et si l'éditeur établit l'absence de code persistant dans la
          passerelle (P10 passe à \NON{})~?
\end{enumerate}
\end{boxexo}

\begin{boxexo}[title={\ding{45}\quad Exercice~16.3. Le U d'un parc \niveaumoyen}]
Reprenez la couverture de SOLITRANS en mai 2025 (Chapitre~\ref{chap:identifier})~:
1\,512 actifs, dont 170 non analysés et 22~\% des serveurs analysés sans
authentification. Proposez une manière de calculer un $U$ du parc sur la
question « l'actif présente-t-il une faille de classe P1~? », et discutez-en
les limites.
\end{boxexo}

\begin{boxexo}[title={\ding{45}\quad Exercice~16.4. Un sujet de recherche \niveauexpert}]
Choisissez l'une des cinq questions ouvertes de la section~16.4 et rédigez
une proposition de sujet de mémoire d'une page~: question, hypothèse,
données, méthode de mesure, critère de validation, limites attendues.
\end{boxexo}

\bigskip

\begin{boxinfo}
\textbf{Pour aller plus loin.}
\begin{itemize}
    \item T.~E.~Minka Mi Nguidjoï, « Vers l'auditologie »~\cite{minka2026auditologie}.
    \item T.~E.~Minka Mi Nguidjoï, sur l'instrument
          trichotomique~\cite{minka2026nek}.
    \item C.~E.~Shannon, « A Mathematical Theory of
          Communication »~\cite{shannon1948}.
\end{itemize}
\end{boxinfo}

\begin{boxinfo}
\textbf{Fin du cours.} La nuit du 14~mars a commencé par vingt-trois
virements et une question~: par où sont-ils entrés~? Elle s'achève sur une
règle d'inventaire, cinq \NEK{} assumés et une passerelle sous scellé. Entre
les deux, l'enquête a préservé, acquis, analysé, daté, remonté à la faille
et rendu compte, en disant à chaque étape ce qu'elle savait et ce qu'elle
ignorait.
\end{boxinfo}
